\section{نتایج محاسباتی و تحلیل}

\subsection{محیط و پروتکل آزمایش}
پیاده‌سازی با \lr{Python 3} انجام شد. \lr{NetworkX} برای انقباض گراف،
\lr{PuLP} و \lr{CBC} برای مدل دقیق، و \lr{NumPy} و \lr{Matplotlib} برای تحلیل
به‌کار رفتند. آزمایش‌ها روی \lr{macOS} با پردازنده \lr{Apple Silicon M4} و
۱۶ گیگابایت حافظه اجرا شدند. برای بازتولید محیط، وابستگی‌ها در فایل
\texttt{requirements.txt} ثبت شده‌اند.

۱۷ نمونه متنی در شش گروه بررسی شدند: ۹ نمونه استاندارد با توزیع یکنواخت و ۸
نمونه مصنوعی دشوار با مخلوطی از اقلام نزدیک نصف ظرفیت، متوسط و کوچک. گروه دوم
فقط از ایده نمونه‌های دشوار ادبیات الهام گرفته است و benchmark رسمی
\lr{Falkenauer} محسوب نمی‌شود. تمام فایل‌ها شامل $n$، ظرفیت، وزن‌ها و یال‌های
چسبندگی‌اند و نام فایل با محتوای آن سازگار است.

برای هر جواب ابتکاری سه شرط به‌صورت برنامه‌ای بررسی شد: هر مؤلفه دقیقاً یک بار
ظاهر شود، ظرفیت هیچ سطلی نقض نشود و پس از بسط مؤلفه‌ها تمام یال‌های چسبندگی
رعایت شوند. هر ۱۷ جواب این اعتبارسنجی را گذراندند.

\subsection{مقدار مرجع و تعریف شکاف}
برای نمونه‌های دارای $n\le30$، مدل دقیق component-level با \lr{CBC} و حد زمان
۱۲۰ ثانیه اجرا شد و فقط وضعیت \lr{Optimal} پذیرفته شد. برای نمونه‌های بزرگ‌تر
از کران پایین معتبر
\begin{equation}
LB=\max\left\{\left\lceil\frac{\sum_c W_c}{C}\right\rceil,
\left|\{c:W_c>C/2\}\right|\right\}
\end{equation}
استفاده شد. شکاف گزارش‌شده برابر فاصله جواب تا مقدار مرجع است. اگر جواب به
$LB$ برسد، بهینگی آن ثابت می‌شود؛ ولی شکاف مثبت نسبت به $LB$ لزوماً برابر
شکاف واقعی تا جواب بهینه نیست. نوع مرجع در همه جدول‌ها صریحاً مشخص شده است.

\subsection{نتایج اصلی}
\begin{table}[htbp]
\centering
\scriptsize
\begin{tabular}{r l r r r r l r r}
\toprule
ردیف & نمونه & $n$ & \lr{FFD} & \lr{Tabu} & مقدار مرجع & مرجع & شکاف (\%) & زمان (s) \\
\midrule
1 & \lr{n=10\_C=100\_pairs=3.txt} & 10 & 2 & 2 & 2 & \lr{ILP} & 0.00 & 0.0000 \\
2 & \lr{n=20\_C=100\_pairs=7.txt} & 20 & 4 & 4 & 4 & \lr{ILP} & 0.00 & 0.0000 \\
3 & \lr{n=30\_C=100\_pairs=10.txt} & 30 & 6 & 6 & 6 & \lr{ILP} & 0.00 & 0.0000 \\
4 & \lr{n=50\_C=200\_pairs=15.txt} & 50 & 9 & 9 & 9 & کران پایین & 0.00 & 0.0000 \\
5 & \lr{n=100\_C=200\_pairs=30.txt} & 100 & 17 & 17 & 17 & کران پایین & 0.00 & 0.0000 \\
6 & \lr{n=150\_C=200\_pairs=45.txt} & 150 & 27 & 27 & 27 & کران پایین & 0.00 & 0.0001 \\
7 & \lr{n=300\_C=500\_pairs=80.txt} & 300 & 51 & 51 & 51 & کران پایین & 0.00 & 0.0002 \\
8 & \lr{n=500\_C=500\_pairs=120.txt} & 500 & 85 & 85 & 85 & کران پایین & 0.00 & 0.0006 \\
9 & \lr{n=1000\_C=1000\_pairs=200.txt} & 1000 & 162 & 162 & 162 & کران پایین & 0.00 & 0.0017 \\
\bottomrule
\end{tabular}
\caption{نتایج نمونه‌های استاندارد؛ نوع مقدار مرجع در ستون «مرجع» مشخص شده است.}
\label{tab:standard-results}
\end{table}

\begin{table}[htbp]
\centering
\scriptsize
\begin{tabular}{r l r r r r l r r}
\toprule
ردیف & نمونه & $n$ & \lr{FFD} & \lr{Tabu} & مقدار مرجع & مرجع & شکاف (\%) & زمان (s) \\
\midrule
1 & \lr{hard\_n=10\_C=100\_pairs=2.txt} & 10 & 4 & 4 & 4 & \lr{ILP} & 0.00 & 0.0001 \\
2 & \lr{hard\_n=20\_C=100\_pairs=5.txt} & 20 & 8 & 8 & 8 & \lr{ILP} & 0.00 & 0.0000 \\
3 & \lr{hard\_n=30\_C=100\_pairs=8.txt} & 30 & 11 & 11 & 11 & \lr{ILP} & 0.00 & 0.0000 \\
4 & \lr{hard\_n=50\_C=200\_pairs=10.txt} & 50 & 17 & 17 & 17 & کران پایین & 0.00 & 0.0000 \\
5 & \lr{hard\_n=100\_C=200\_pairs=20.txt} & 100 & 33 & 33 & 33 & کران پایین & 0.00 & 0.0001 \\
6 & \lr{hard\_n=150\_C=200\_pairs=30.txt} & 150 & 50 & 50 & 49 & کران پایین & 2.04 & 0.1152 \\
7 & \lr{hard\_n=300\_C=500\_pairs=50.txt} & 300 & 100 & 99 & 99 & کران پایین & 0.00 & 0.0130 \\
8 & \lr{hard\_n=500\_C=500\_pairs=80.txt} & 500 & 167 & 166 & 166 & کران پایین & 0.00 & 0.3166 \\
\bottomrule
\end{tabular}
\caption{نتایج نمونه‌های مصنوعی دشوار؛ نوع مقدار مرجع در ستون «مرجع» مشخص شده است.}
\label{tab:hard-results}
\end{table}


در هر شش نمونه کوچک استاندارد و دشوار، جواب \lr{FFD+Tabu} با جواب بهینه مدل
دقیق برابر شد. هر ۹ نمونه استاندارد نیز به کران پایین رسیدند. این نتیجه بیش از
آنکه نشانه دشواری جستجو باشد، نشان می‌دهد \lr{FFD} برای این مجموعه یکنواخت
نقطه شروع بسیار قوی‌ای است؛ در نتیجه الگوریتم با تشخیص رسیدن به مرجع معتبر
بلافاصله متوقف می‌شود.

در نمونه مصنوعی دشوار ۱۵۰ قلمی، جواب ۵۰ سطل و کران پایین ۴۹ است؛ بنابراین
فاصله نسبت به مرجع ۲٫۰۴ درصد باقی ماند. در مقابل، روی نمونه‌های دشوار ۳۰۰ و
۵۰۰ قلمی، \lr{FFD} به‌ترتیب ۱۰۰ و ۱۶۷ سطل ساخت و جستجوی تابو هر دو را یک سطل
بهبود داد و به کران‌های پایین ۹۹ و ۱۶۶ رسید. شکل~\ref{fig:hard-improvement}
این تفاوت را نشان می‌دهد.

\begin{figure}[htbp]
\centering
\includegraphics[width=0.78\textwidth]{figures/fig6_pathological.pdf}
\caption{مقایسه \lr{FFD}، جواب نهایی تابو و کران پایین روی سه نمونه مصنوعی دشوار.}
\label{fig:hard-improvement}
\end{figure}

\subsection{تحلیل پایداری با ۳۰ شروع مستقل}
برای هر نمونه متوسط استاندارد و دشوار، الگوریتم ۳۰ بار با شروع تصادفی و seedهای
۱ تا ۳۰ اجرا شد. جدول~\ref{tab:robustness} هم میانگین جواب اولیه و هم جواب
نهایی را گزارش می‌کند؛ بنابراین اثر خود جستجوی تابو از کیفیت \lr{FFD} تفکیک
می‌شود.

\begin{table}[htbp]
\centering
\scriptsize
\begin{tabular}{lrrrrrr}
\toprule
نمونه & اجرا & میانگین اولیه & میانگین نهایی & انحراف معیار نهایی & بهترین & زمان میانگین (s) \\
\midrule
\lr{Std n=50} & 30 & 9.07 & 9.00 & 0.00 & 9 & 0.0001 \\
\lr{Std n=100} & 30 & 17.07 & 17.00 & 0.00 & 17 & 0.0001 \\
\lr{Std n=150} & 30 & 28.20 & 27.00 & 0.00 & 27 & 0.0130 \\
\lr{Hard n=50} & 30 & 18.33 & 17.00 & 0.00 & 17 & 0.0014 \\
\lr{Hard n=100} & 30 & 36.20 & 33.43 & 0.50 & 33 & 0.0459 \\
\lr{Hard n=150} & 30 & 53.53 & 50.00 & 0.00 & 50 & 0.1703 \\
\bottomrule
\end{tabular}
\caption{تحلیل پایداری در ۳۰ شروع تصادفی مستقل؛ اعداد مستقیماً از فایل \lr{robustness.csv} تولید شده‌اند.}
\label{tab:robustness}
\end{table}


در نمونه استاندارد ۱۵۰ قلمی، میانگین تعداد سطل از ۲۸٫۲۰ به ۲۷ کاهش یافت. در
نمونه‌های دشوار ۵۰، ۱۰۰ و ۱۵۰ قلمی، کاهش میانگین به‌ترتیب از ۱۸٫۳۳ به ۱۷، از
۳۶٫۲۰ به ۳۳٫۴۳ و از ۵۳٫۵۳ به ۵۰ بود. انحراف معیار نهایی نمونه دشوار ۱۰۰ قلمی
۰٫۵۰ است؛ این مورد نشان می‌دهد نقطه شروع هنوز بر کیفیت نهایی اثر محدود دارد.
سایر نمونه‌ها در ۳۰ اجرا به یک مقدار نهایی مشترک رسیدند.

\begin{figure}[htbp]
\centering
\includegraphics[width=0.88\textwidth]{figures/fig4_boxplot.pdf}
\caption{توزیع جواب نهایی در ۳۰ شروع تصادفی مستقل.}
\label{fig:boxplot}
\end{figure}

\begin{figure}[htbp]
\centering
\includegraphics[width=0.9\textwidth]{figures/gap_comparison_random_vs_ffd.png}
\caption{فاصله میانگین شروع تصادفی و جواب پس از تابو نسبت به کران پایین.}
\label{fig:random-gap}
\end{figure}

\begin{figure}[htbp]
\centering
\includegraphics[width=0.9\textwidth]{figures/tabu_improvement_on_random.png}
\caption{میانگین تعداد سطل حذف‌شده توسط جستجوی تابو در ۳۰ اجرا.}
\label{fig:mean-improvement}
\end{figure}

\subsection{همگرایی، زمان و ساختار جواب}
شکل~\ref{fig:convergence} بهترین جواب مشاهده‌شده را در طول جستجو برای سه شروع
تصادفی نشان می‌دهد. شکل~\ref{fig:time} زمان wall-clock مدل دقیق و روش ابتکاری
را برای نمونه‌های کوچک روی مقیاس لگاریتمی مقایسه می‌کند. زمان‌های بسیار کوچک
ابتکاری شامل contraction، ساخت جواب اولیه و اعتبارسنجی‌اند و در صورت رسیدن
جواب اولیه به مقدار مرجع، حلقه جستجو اجرا نمی‌شود.

\begin{figure}[htbp]
\centering
\includegraphics[width=0.85\textwidth]{figures/fig1_convergence.pdf}
\caption{روند بهترین جواب در جستجوی تابو از شروع‌های تصادفی واقعی.}
\label{fig:convergence}
\end{figure}

\begin{figure}[htbp]
\centering
\includegraphics[width=0.75\textwidth]{figures/fig2_time_comparison.pdf}
\caption{مقایسه زمان مدل دقیق و روش ابتکاری در نمونه‌های کوچک؛ محور عمودی لگاریتمی است.}
\label{fig:time}
\end{figure}

شکل~\ref{fig:gap-size} فاصله تا مرجع را بر حسب اندازه و
شکل~\ref{fig:packing} یک جواب component-level اعتبارسنجی‌شده را نمایش می‌دهد.

\begin{figure}[htbp]
\centering
\begin{minipage}{0.48\textwidth}
\centering
\includegraphics[width=\linewidth]{figures/fig3_gap_vs_size.pdf}
\caption{فاصله جواب تا مرجع اعلام‌شده.}
\label{fig:gap-size}
\end{minipage}\hfill
\begin{minipage}{0.48\textwidth}
\centering
\includegraphics[width=\linewidth]{figures/fig5_bin_visualization.pdf}
\caption{چیدمان مؤلفه‌ها در نمونه ۳۰ قلمی.}
\label{fig:packing}
\end{minipage}
\end{figure}

\subsection{تهدیدهای اعتبار}
نمونه‌ها مصنوعی و تعداد آن‌ها محدود است؛ بنابراین نتایج قابل تعمیم مستقیم به
همه benchmarkهای \lr{BPP} نیست. همچنین کران پایین استفاده‌شده برای نمونه‌های
بزرگ ساده است و در حالت فاصله مثبت، مقدار بهینه ناشناخته می‌ماند. زمان‌های
اجرا فقط روی یک سامانه اندازه‌گیری شده‌اند. بااین‌حال، ثبت seedها، انتشار تمام
فایل‌های داده، تولید خودکار جدول‌ها و شکل‌ها و اعتبارسنجی مستقل جواب‌ها امکان
بازتولید و بررسی نتایج حاضر را فراهم می‌کند.
